\documentclass[11pt,a4paper]{article}
\usepackage[utf8]{inputenc}
\usepackage[T1]{fontenc}
\usepackage[margin=0.72in]{geometry}
\usepackage{enumitem}
\usepackage{titlesec}
\usepackage{xcolor}
\usepackage{hyperref}
\definecolor{primarycolor}{RGB}{33,37,41}
\definecolor{secondarycolor}{RGB}{0,92,153}
\definecolor{accentcolor}{RGB}{95,105,115}
\hypersetup{colorlinks=true,urlcolor=secondarycolor,linkcolor=secondarycolor}
\titleformat{\section}{\large\bfseries\color{primarycolor}}{}{0em}{}
\titlespacing*{\section}{0pt}{0.9em}{0.35em}
\setlist[itemize]{leftmargin=1.2em,itemsep=0.12em,topsep=0.2em}
\pagestyle{empty}
\begin{document}
\begin{center}
{\LARGE\bfseries\color{primarycolor} Lena Varga-Holm}\\
{\large Senior Data Engineer | Data Platform}\\[0.25em]
Email: applicant.inbox@example.com | Phone: +44 7700 900123
\end{center}

\section*{PROFESSIONAL SUMMARY}
Senior Data Engineer with experience owning ingestion and orchestration for a data platform built on Airflow, Snowflake, dbt and Kafka. Designed and operated production data pipelines processing around 30 million events per day, with end-to-end latency under two minutes at p95. Delivered measurable improvements in warehouse build time, database load and cloud warehouse costs. Experienced in Python, SQL, data quality, open-source contributions and mentoring engineers.

\section*{EXPERIENCE}
\textbf{Senior Data Engineer | Lukovac Freight Hub d.o.o. | 2022--Present}
\begin{itemize}
\item Own ingestion and orchestration for the data platform using Airflow, Snowflake, dbt and Kafka; participate in the data-platform on-call rotation one week in six.
\item Designed and built a shipment-event streaming pipeline using partner webhooks, Kafka, Kafka Connect and Snowpipe Streaming, processing around 30 million events per day with end-to-end latency under two minutes at p95.
\item Added a dead-letter topic and replay tooling so failed events can be reprocessed without data loss.
\item Migrated approximately 300 stored procedures into dbt models with tests and documentation; introduced incremental models and slim CI using state:modified in GitHub Actions. Cut the nightly warehouse build from 5 hours 10 minutes to 1 hour 35 minutes and made lineage visible to analysts.
\item Replaced hourly full-table extracts from the orders database with Debezium change data capture into Kafka, reducing replica CPU during loads from around 70\% to around 25\%.
\item Reduced the Snowflake bill from approximately \$38,000 to \$24,000 per month in 2024 through warehouse sizing, auto-suspend and removing unused tables.
\item Mentor two mid-level engineers and run data platform office hours.
\end{itemize}

\textbf{Data Engineer | Kovalent Pay | 2019--2022}
\begin{itemize}
\item Built PySpark batch jobs on AWS EMR to match bank settlement files against internal ledgers for a transaction reconciliation engine.
\item Reduced the finance team's manual reconciliation work from around three days to around four hours per month-end.
\item Worked on Redshift warehouse design and data quality using Great Expectations.
\end{itemize}

\section*{OPEN-SOURCE CONTRIBUTIONS}
\begin{itemize}
\item Fixed a pagination bug and added tests for a community Airflow provider package for a SaaS CRM API; two pull requests were merged. Technologies: Python, Airflow, pytest.
\end{itemize}

\section*{SKILLS}
Python, SQL, PySpark, Airflow, dbt, Kafka, Kafka Connect, Debezium, Snowflake, AWS MSK, AWS EMR, S3, Glue Data Catalog, Redshift, PostgreSQL, Terraform, GitHub Actions, Great Expectations, pytest, data ingestion and orchestration, streaming pipelines, change data capture, data quality, warehouse optimisation.

\section*{CERTIFICATION}
AWS Certified Data Engineer -- Associate | May 2024

\section*{EDUCATION}
MSc in Computer Science | Northgate Technical University\\
Thesis: Incremental materialized views
\end{document}